\documentclass[11pt]{article}
\usepackage[margin=1in]{geometry}
\usepackage{amsmath,amssymb,amsthm}
\usepackage{physics}
\usepackage{enumitem}
\usepackage{hyperref}
\usepackage{titlesec}

\theoremstyle{definition}
\newtheorem{definition}{Definition}[section]
\newtheorem{example}{Example}[section]
\newtheorem{remark}{Remark}[section]
\newtheorem{proposition}{Proposition}[section]
\newtheorem{theorem}{Theorem}[section]

\titleformat{\section}{\large\bfseries}{Lecture 14.\thesection}{1em}{}

\title{\textbf{Lecture 14: The Stabilizer Formalism} \\ \large Understanding Quantum Information \& Computation}
\author{Based on the lectures by John Watrous (IBM Quantum)}
\date{}

\begin{document}
\maketitle
\tableofcontents
\newpage

%----------------------------------------------------------------------
\section{Overview}
%----------------------------------------------------------------------

The \textbf{stabilizer formalism} is a powerful algebraic framework for describing and analysing quantum error-correcting codes.  It provides:
\begin{enumerate}
    \item A compact way to specify quantum codes via groups of Pauli operators.
    \item Systematic methods for syndrome extraction and error correction.
    \item A bridge to the theory of classical linear codes.
\end{enumerate}

%----------------------------------------------------------------------
\section{The Pauli Group}
%----------------------------------------------------------------------

\subsection{Single-Qubit Pauli Group}

\begin{definition}
The \textbf{single-qubit Pauli group} $\mathcal{P}_1$ consists of the 16 operators
\[
    \mathcal{P}_1 = \{\pm I,\;\pm iI,\;\pm X,\;\pm iX,\;\pm Y,\;\pm iY,\;\pm Z,\;\pm iZ\}.
\]
These form a group under matrix multiplication.
\end{definition}

Key properties of Pauli matrices:
\begin{itemize}
    \item $X^2=Y^2=Z^2=I$.
    \item $XY=iZ$, $YZ=iX$, $ZX=iY$ (cyclic).
    \item Any two distinct non-identity Paulis either commute or anticommute.
\end{itemize}

\subsection{$n$-Qubit Pauli Group}

\begin{definition}
The \textbf{$n$-qubit Pauli group} $\mathcal{P}_n$ consists of all $n$-fold tensor products of single-qubit Pauli operators, multiplied by overall phases $\{\pm 1,\pm i\}$:
\[
    \mathcal{P}_n = \bigl\{\alpha\,\sigma_1\otimes\cdots\otimes\sigma_n : \alpha\in\{1,-1,i,-i\},\;\sigma_k\in\{I,X,Y,Z\}\bigr\}.
\]
$|\mathcal{P}_n|=4^{n+1}$.
\end{definition}

\begin{example}
$n=3$: $X\otimes I\otimes Z$, $\;-Y\otimes Y\otimes I$, $\;iI\otimes X\otimes X$, etc.
\end{example}

\paragraph{Commutation.}  Two elements of $\mathcal{P}_n$ either commute or anticommute.  Whether they commute can be determined by counting the number of positions where the corresponding single-qubit Paulis anticommute: commute iff this count is even.

%----------------------------------------------------------------------
\section{Stabilizer Groups and Stabilizer States}
%----------------------------------------------------------------------

\begin{definition}[Stabilizer group]
A \textbf{stabilizer group} $\mathcal{S}$ is an abelian subgroup of $\mathcal{P}_n$ such that $-I\notin\mathcal{S}$.
\end{definition}

\begin{definition}[Stabilizer state / code space]
The \textbf{stabilizer space} (or \textbf{code space}) of $\mathcal{S}$ is
\[
    V_{\mathcal{S}} = \bigl\{\ket{\psi} : g\ket{\psi}=\ket{\psi}\;\;\forall\,g\in\mathcal{S}\bigr\}.
\]
Every vector in $V_{\mathcal{S}}$ is a $+1$ eigenvector of every element of $\mathcal{S}$.
\end{definition}

\begin{proposition}
If $\mathcal{S}$ is generated by $n-k$ independent generators $g_1,\dots,g_{n-k}$ (each a Hermitian element of $\mathcal{P}_n$), then $\dim V_{\mathcal{S}}=2^k$.  The code encodes $k$ \textbf{logical qubits} into $n$ \textbf{physical qubits}.
\end{proposition}

\paragraph{Generators.}  A stabilizer group with $n-k$ generators is compactly specified by listing just those generators:
\[
    \mathcal{S} = \langle g_1,g_2,\dots,g_{n-k}\rangle.
\]

\begin{example}[3-qubit bit-flip code]
$n=3$, $k=1$.  Generators:
\[
    g_1 = Z\otimes Z\otimes I,\qquad g_2 = I\otimes Z\otimes Z.
\]
Code space: $\operatorname{span}\{\ket{000},\ket{111}\}$, i.e.\ $\ket{0}_L=\ket{000}$, $\ket{1}_L=\ket{111}$.
\end{example}

\begin{example}[3-qubit phase-flip code]
$n=3$, $k=1$.  Generators:
\[
    g_1 = X\otimes X\otimes I,\qquad g_2 = I\otimes X\otimes X.
\]
Code space: $\operatorname{span}\{\ket{+++},\ket{---}\}$.
\end{example}

%----------------------------------------------------------------------
\section{Error Detection and Correction with Stabilizers}
%----------------------------------------------------------------------

\subsection{Syndrome Extraction}

For an error $E\in\mathcal{P}_n$ acting on a codeword $\ket{\psi}\in V_{\mathcal{S}}$, the \textbf{syndrome} is the vector of $\pm 1$ eigenvalues obtained by measuring each generator:
\[
    g_j\,(E\ket{\psi}) = \pm\,E\ket{\psi}.
\]
Specifically, $g_j(E\ket{\psi})=+E\ket{\psi}$ if $g_j$ and $E$ commute, and $g_j(E\ket{\psi})=-E\ket{\psi}$ if they anticommute.

\begin{definition}[Syndrome]
The syndrome of error $E$ with respect to generators $g_1,\dots,g_{n-k}$ is the binary string
\[
    s(E) = (s_1,\dots,s_{n-k}),\qquad s_j = \begin{cases}0 & [g_j,E]=0,\\ 1 & \{g_j,E\}=0.\end{cases}
\]
\end{definition}

\subsection{Correctable Errors}

\begin{theorem}
A stabilizer code with stabilizer group $\mathcal{S}$ can correct a set of errors $\{E_1,\dots,E_r\}$ if and only if for every pair $E_i,E_j$:
\[
    E_i^\dagger E_j \notin \mathcal{N}(\mathcal{S})\setminus\mathcal{S},
\]
where $\mathcal{N}(\mathcal{S})$ is the \textbf{normaliser} of $\mathcal{S}$ in $\mathcal{P}_n$ (the set of Pauli operators that commute with every element of $\mathcal{S}$).
\end{theorem}

Equivalently: distinct errors must have distinct syndromes, or their product must lie in $\mathcal{S}$ (so they have the same effect on the code space and no correction is needed).

\subsection{Correction Procedure}

\begin{enumerate}
    \item Measure the generators $g_1,\dots,g_{n-k}$ to obtain the syndrome $s$.
    \item Look up the error $E$ associated with syndrome $s$ (from a pre-computed table).
    \item Apply $E^\dagger$ (or equivalently $E$, since Pauli operators are Hermitian and unitary) to correct the error.
\end{enumerate}

%----------------------------------------------------------------------
\section{The Normaliser and Logical Operators}
%----------------------------------------------------------------------

\begin{definition}
The \textbf{normaliser} $\mathcal{N}(\mathcal{S})$ consists of all elements of $\mathcal{P}_n$ that commute with every element of $\mathcal{S}$.  Elements of $\mathcal{N}(\mathcal{S})\setminus\mathcal{S}$ act non-trivially on the code space and represent \textbf{logical operations} on the encoded qubits.
\end{definition}

For an $\llbracket n,k\rrbracket$ code, one can find $2k$ independent logical operators $\bar{X}_1,\bar{Z}_1,\dots,\bar{X}_k,\bar{Z}_k$ in $\mathcal{N}(\mathcal{S})\setminus\mathcal{S}$ satisfying the same commutation relations as Pauli operators on $k$ qubits.

%----------------------------------------------------------------------
\section{The $\llbracket n,k,d\rrbracket$ Notation}
%----------------------------------------------------------------------

\begin{definition}
An $\llbracket n,k,d\rrbracket$ \textbf{stabilizer code} encodes $k$ logical qubits into $n$ physical qubits with \textbf{distance} $d$, defined as the minimum weight of any element in $\mathcal{N}(\mathcal{S})\setminus\mathcal{S}$.

The code can detect up to $d-1$ errors and correct up to $\lfloor(d-1)/2\rfloor$ errors.
\end{definition}

\begin{example}\leavevmode
\begin{itemize}
    \item 3-qubit bit-flip code: $\llbracket 3,1,1\rrbracket$ (distance 1 for general errors, but distance 3 restricted to $X$ errors).
    \item Shor code: $\llbracket 9,1,3\rrbracket$ --- corrects any single-qubit error.
    \item Steane code: $\llbracket 7,1,3\rrbracket$ --- more efficient than Shor's code.
    \item 5-qubit code: $\llbracket 5,1,3\rrbracket$ --- the smallest code correcting any single-qubit error.
\end{itemize}
\end{example}

%----------------------------------------------------------------------
\section{Connection to Classical Codes}
%----------------------------------------------------------------------

Many stabilizer codes are constructed from classical linear codes via the \textbf{CSS (Calderbank--Shor--Steane) construction}:

\begin{theorem}[CSS codes]
Given two classical linear codes $C_1$ and $C_2$ with $C_2\subset C_1$, one can construct a quantum stabilizer code $\llbracket n,k,d\rrbracket$ where $k=\dim(C_1)-\dim(C_2)$ and $d$ depends on the minimum distances of $C_1$ and $C_2^\perp$.
\end{theorem}

The Steane code is a CSS code constructed from the classical $[7,4,3]$ Hamming code.

%----------------------------------------------------------------------
\section{Summary}
%----------------------------------------------------------------------

\begin{itemize}
    \item The \textbf{Pauli group} $\mathcal{P}_n$ provides the algebraic foundation for stabilizer codes.
    \item A \textbf{stabilizer code} is specified by an abelian subgroup $\mathcal{S}\subset\mathcal{P}_n$ (not containing $-I$); the code space consists of the common $+1$ eigenspace.
    \item \textbf{Syndrome measurement} of the generators identifies errors without disturbing the encoded information.
    \item The \textbf{normaliser} $\mathcal{N}(\mathcal{S})$ determines logical operators and the code distance.
    \item The $\llbracket n,k,d\rrbracket$ notation compactly describes a code's parameters.
    \item \textbf{CSS codes} bridge classical and quantum coding theory.
\end{itemize}

\end{document}
